% ---- pg-515 （书页 503） ----

的训练集来确定的常数(参数).

网络结构是计算单元的组织形式及允许的连接类型. 在统计应用中, 计算单元被安排成若干层, 而且不同层的结点之间是连接的, 但在同一层中的结点不连接. 接受原始输入的层称为输入层, 最后一层称为输出层, 位于输入层和输出层之间的层统称为隐层. \hyperref[fig:11.20]{图 11.20} 是一个多层神经网络的简单示意图.

\begin{mvfigure}
\begin{tikzpicture}[scale=0.9,transform shape,
  nd/.style={circle,draw,minimum size=10mm,inner sep=0pt}]
% 输入层
\foreach \i/\x in {1/0,2/1.6,3/3.2,4/4.8}{
  \node[nd] (I\i) at (\x,0) {};
  \draw[-latex] (\x,-1.05)--(\x,-0.52);
  \node[below] at (\x,-1.0) {$X_{\i}$};
}
% 中间层(隐层)
\node[nd] (H1) at (1.6,1.9) {};
\node[nd] (H2) at (3.2,1.9) {};
% 输出层
\node[nd] (O1) at (0.8,3.8) {};
\node[nd] (O2) at (2.4,3.8) {};
\node[nd] (O3) at (4.0,3.8) {};
% 连接
\foreach \i in {1,...,4}{
  \draw (I\i)--(H1); \draw (I\i)--(H2);
}
\foreach \h in {1,2}{
  \draw (H\h)--(O1); \draw (H\h)--(O2); \draw (H\h)--(O3);
}
% 输出箭头
\foreach \x in {0.8,2.4,4.0}{
  \draw[-latex] (\x,4.35)--(\x,4.95);
}
% 层名
\node[anchor=east] at (-0.85,3.8) {输出层};
\node[anchor=east] at (-0.85,1.9) {中间层(隐层)};
\node[anchor=east] at (-0.85,0) {输入层};
\end{tikzpicture}
\phantomsection\label{fig:11.20}\mvcap{图 11.20\quad 有一个隐层的神经网络}
\end{mvfigure}

神经网络可用于判别和分类. 在这种应用中, 输入变量为所测量的总体特性量 $X_1,X_2,\cdots,X_p$, 而输出变量则为指示各组归属的类别变量. 最新的实践经验表明, 正确构造的神经网络其表现与本章中所讨论的逻辑斯谛回归和判别函数一样好. 文献[26]对神经网络在应用统计学中的运用提供了很好的讨论.

\mvsub{变量选择}

在判别分析的某些应用中, 可以获得大量变量的数据. 穆西阿迪和戈斯((Mucciardi \& Gose, 见文献[27])讨论了基于 157 个变量的判别分析.\footnote{想象一下验证关于 157 元正态性的假定和参数的联合估计问题, 或设想一下 $157\times 157$ 公共协方差阵中的 12\,403 个参数!} 在这种情况下, 显然希望选择一个包含较少变量, 同时又能保留与原集合几乎一样多信息的子集, 这就是逐步判别分析法的目标, 几个流行的商用计算机程序具有这种功能.

如果使用了逐步判别分析法(或任何一种变量选择方法), 则在解释其分析结果时必须小心(见文献[28]). 不管使用何种准则来选择变量, 所得到的变量子集不一定是最优的. 例如, 根据最小化表观误差率或最大化“判别能力”所得到的子集在用于未来样本时可能效果很差. 如果在各变量之间或变量的各线性组合之间存在强相关, 则与变量选择方法相联系的问题会变得更严重.

对给定的一个数据集, 若目标是分类, 则选择一个最优变量子集的问题变得特别麻烦. 至少, 导出的分类函数应该用一个有效的样本来评估. 正如默里(Murray, 见文献[28])所建议的那样, 将样本分成若干组并对每一组确定“最优”变量子集, 也许是较好的想法. 对一个给定的变量, 它在这些最优子集中出现的次数就是这个变量在未来分类中的价值度量.

\mvsub{检验组间的差别}

在两组数据的分类问题中我们已经指出, 除非两个总体分离得很好, 否则有效分配大概是不可能的. 对于两组以上数据, 情况也一样, 除非总体的均值向量之间有显著差别, 否则就不要做分类. 假定数据的分布接近多元正态分布, 且共有相同的协方差阵, 此时可用 MANOVA (多元方差分析)来检验总体均值向量之间的差别. 虽然表观的显著差别并不自动意味着有效的分类, 但检验是必要的第一步. 如果没有发现(总体均值的)显著差别, 则构造分类法则可能只是浪费时间.

\mvsub{图}

先进的计算机作图使人们现在可以用目测方法直观地考察二维或三维的多元数据. 这样,
